\section{Deneysel Kurulum}
\label{sec:setup}

\subsection{Veri Kümesi}

UCI HAR~\cite{anguita2013public} veri kümesinin bir uzantısı olan
Human Activities and Postural Transitions (\hapt{})
veri kümesini~\cite{reyes2015transition} kullanıyoruz. \hapt{},
belde taşınan bir akıllı telefondan 50~Hz'de kaydedilmiş üç eksenli
doğrusal ivme verilerini içerir; 30 katılımcı altı temel aktiviteyi
gerçekleştirir: \textsc{walking}, \textsc{upstairs},
\textsc{downstairs}, \textsc{sitting}, \textsc{standing} ve
\textsc{laying}. \cite{reyes2015transition} çalışmasındaki kanonik
özne-ayrık eğitim/doğrulama/test bölmesini benimsiyoruz. Bu bölme,
uzunluğu 128 olan (50~Hz'de 2.56~s) 7{,}352 eğitim penceresi,
1{,}515 doğrulama penceresi ve 3{,}399 test penceresi üretir. Bu
makalede raporlanan tüm doğruluk ve F1 değerleri ayrılmış test
bölmesi üzerinde hesaplanmıştır.

Akış değerlendirmesinde her test penceresini 50~Hz'de örnek örnek
besliyor ve her pencere sınırında gizli durumu sıfırlıyoruz. Böylece
değerlendirme, gömülü çalışma zamanındaki davranışla birebir eşleşir.

\subsection{Eğitim Protokolü}

Modeller, masaüstü CPU üzerinde PyTorch~2.x ve Adam iyileştiricisi
($\eta{=}10^{-3}$, yığın boyutu~64, 100 epoch) ile eğitilmiştir. IHT
seyrekleştirme maskesi Bölüm~\ref{sec:method:iht}'deki kübik
çizelgeyi izler; hedef seyreklik epoch~50'de ulaşılır ve sonraki
50 epoch'luk ince ayar boyunca maske sabit kalır. Raporlanan her
yapılandırma beş tohum $\{0, 1, 2, 3, 4\}$ üzerinde tekrarlanır;
aksi belirtilmedikçe istatistikler ortalama~$\pm$~standart sapma
olarak verilmiştir.

\subsection{Gömülü Hedef Donanım}

İki bare-metal hedef değerlendirilmiştir:
\begin{itemize}
\item \textbf{\arduino{}} (ATmega328P): 16~MHz'de çalışan 8-bit AVR
çekirdeği; 32~KB Flash; 2~KB SRAM; donanım $8{\times}8$ çarpıcısı;
kayan nokta birimi \emph{yok}. Araç zinciri: \texttt{-Os}
seçeneğiyle \texttt{avr-gcc} kullanan Arduino~IDE~2.3.x.
\item \textbf{\msp{}}: kalibre edilmiş 1~MHz DCO'da çalışan 16-bit
MSP430 çekirdeği; 16~KB Flash; 512~B SRAM; \emph{hiçbir tür donanım
çarpıcısı yok}; FPU yok. Araç zinciri: \texttt{-O2} ile derlenen TI
\texttt{msp430-elf-gcc} bare-metal derlemesini kullanan Code
Composer Studio~12.x. Energia çalışma zamanı \emph{kullanılmamıştır}.
\end{itemize}

Her iki hedefte de yüklenen program şu bileşenlerden oluşur:
(i)~\texttt{fastgrnn.cpp} çıkarım motoru,
(ii)~Q15 ağırlık tablosu ve tensör-başına ölçekler
(\texttt{model\_weights.h}),
(iii)~256 girişli sigmoid ve tanh arama tabloları ve
(iv)~9600 (\msp{}) ya da 115{,}200~baud (\arduino{}) hızında sınıf
etiketi gönderen UART sürücüsü. Geri kalan Flash alanı kullanılmaz.

\subsection{Platformlar Arası Doğrulama Protokolü}
\label{sec:setup:verify}

FP32 PyTorch referansı ile Q15 C uygulaması arasında \%100 tahmin
uyumu iddiasını doğrulamak için (Bölüm~\ref{sec:results:crossplatform}),
aynı girdileri iki çıkarım yolundan geçirip sonuçları
karşılaştırıyoruz. Bir Python denetim betiği
(\texttt{test\_inference\_python.py}) gömülü hedefe aktarılacak Q15
ağırlıkları yükler, NumPy içinde C ile eşdeğer bir ileri geçiş
çalıştırır (LUT aktivasyonları ve tensör-başına ölçek çözme adımları
dahil) ve her test penceresi için \texttt{argmax} tahminlerini FP32
referansıyla karşılaştırır. Arduino ve MSP430 üzerindeki C uygulaması
da pencere-başına tahminleri UART üzerinden gönderir; bu tahminler
aynı referansla denetlenir. Üç yürütme yolu -- FP32 PyTorch,
NumPy'deki C-eşdeğer uygulama ve bare-metal C -- 3{,}399 pencerenin
tamamında aynı sınıfı üretir.

\subsection{Canlı Sensör Yapılandırması}

Canlı gösterimler için InvenSense MPU-6050 sunan bir GY-521 modülü
I\textsuperscript{2}C üzerinden bağlanır. \msp{} üzerinde USCI\_B0
doğrudan sürülür (P1.6~$=$~SCL, P1.7~$=$~SDA, Energia
\texttt{Wire} katmanı yoktur); \arduino{} üzerinde standart Wire
kütüphanesi kullanılır. FP32 eğitim verisi ölçeklemesiyle eşleşmesi
için yalnızca $\pm 2$~g aralığındaki ivmeölçer okunur. Yazılım
zamanlayıcısı örnekleme döngüsünü 50~Hz'e ayarlar; çıkarım
adımından sonra örnek başına 7--11~ms zaman payı kalır
(Bölüm~\ref{sec:results:streaming}).
